% ============================================================================
% article_vs_koma.tex - KOMA-Script Document Class (scrartcl)
% ============================================================================
% Demonstrates the KOMA-Script replacement for the standard article class,
% focusing on parskip, heading styles, and the DIV option for typearea.
%
% Compilation:
%   pdflatex article_vs_koma.tex
%
% Requires: KOMA-Script (koma-script package in TeX Live / MiKTeX)
% ============================================================================

% --- scrartcl replaces article; compatible but with KOMA-Script extensions ---
\documentclass[
    12pt,               % base font size
    a4paper,            % paper size
    parskip=half,       % paragraph separation: half-line space, no indent
    headings=small,     % smaller heading fonts (options: big, normal, small)
    DIV=12              % typearea calculation: higher = narrower margins
]{scrartcl}

% ============================================================================
% PACKAGE IMPORTS
% ============================================================================

\usepackage[utf8]{inputenc}
\usepackage[T1]{fontenc}
\usepackage{lmodern}
\usepackage[english]{babel}
\usepackage{microtype}

\usepackage{xcolor}
\usepackage{booktabs}   % Professional tables
\usepackage{enumitem}   % List customisation

\usepackage[
    backend=biber,
    style=numeric
]{biblatex}

\usepackage{hyperref}
\hypersetup{
    colorlinks=true,
    linkcolor=blue,
    citecolor=teal,
    urlcolor=cyan,
    pdftitle={KOMA-Script scrartcl Demo},
    pdfauthor={LaTeX Student}
}

% ============================================================================
% KOMA-SCRIPT SPECIFIC SETTINGS
% ============================================================================

% Subject line (KOMA-Script only) — appears between title and author
\subject{KOMA-Script Class Demonstration}

% Subtitle (KOMA-Script only)
\subtitle{A comparison with the standard \texttt{article} class}

% ============================================================================
% DOCUMENT METADATA
% ============================================================================

\title{KOMA-Script: \texttt{scrartcl} in Practice}
\author{LaTeX Student}
\date{\today}

% ============================================================================
% DOCUMENT BODY
% ============================================================================

\begin{document}

\maketitle

% ============================================================================
% ABSTRACT — KOMA-Script uses the abstract environment too,
% but styles it differently from the standard class.
% ============================================================================

\begin{abstract}
KOMA-Script is a mature, European-oriented replacement for the standard
\LaTeX{} document classes. This document demonstrates the \texttt{scrartcl}
class, which replaces \texttt{article}, and highlights the three most
practically significant options: \texttt{parskip} for paragraph spacing,
\texttt{headings} for heading size, and \texttt{DIV} for the typearea
(text block) calculation.
\end{abstract}

\tableofcontents

% ============================================================================
% SECTION 1: MOTIVATION
% ============================================================================

\section{Why KOMA-Script?}
\label{sec:motivation}

The standard \LaTeX{} document classes (\texttt{article}, \texttt{report},
\texttt{book}) were designed in the 1980s, primarily for American typographic
conventions. KOMA-Script re-implements all three classes—as \texttt{scrartcl},
\texttt{scrreprt}, and \texttt{scrbook}—with the following improvements:

\begin{itemize}
    \item Automatic typearea calculation based on classical typographic ratios
          (the \texttt{DIV} option or the \texttt{typearea} package)
    \item Configurable paragraph spacing (\texttt{parskip} option)
    \item Flexible heading styles without rewriting \verb|\@makechapterhead|
    \item Built-in support for scrpage headers, KOMA options for captions,
          and an extensive scripting interface via \texttt{KOMA-Options}
    \item Better defaults for European paper sizes (A4, A5, B5)
\end{itemize}

\subsection{Drop-In Replacement}

For most documents, switching from \texttt{article} to \texttt{scrartcl}
requires only changing the \verb|\documentclass| line. Existing commands
such as \verb|\section|, \verb|\subsection|, \verb|\begin{abstract}|, and
\verb|\maketitle| work unchanged.

% ============================================================================
% SECTION 2: THE parskip OPTION
% ============================================================================

\section{Paragraph Separation: the \texttt{parskip} Option}
\label{sec:parskip}

By default, standard \LaTeX{} separates paragraphs with an indent and no
vertical space. KOMA-Script's \texttt{parskip} option switches to the
European convention of a blank half-line between paragraphs.

\subsection{Available Values}

\begin{description}
    \item[\texttt{parskip=false}]
        Default \LaTeX{} behaviour: first-line indent, no extra vertical space.
    \item[\texttt{parskip=true} or \texttt{full}]
        One full line of space between paragraphs; no indent on any paragraph.
    \item[\texttt{parskip=half}]
        Half a line of space between paragraphs; no indent. (\emph{Used in
        this document.})
    \item[\texttt{parskip=full*} / \texttt{half*}]
        Same as above, but the last line of a paragraph may be short without
        triggering an overfull box warning.
\end{description}

In this document the option \texttt{parskip=half} is active. Notice that
paragraphs are separated by a small gap rather than an indentation.

This is the second paragraph. It follows the first without any first-line
indent, but with a half-line vertical gap. This style is common in German
and other European academic and commercial typesetting traditions.

% ============================================================================
% SECTION 3: THE headings OPTION
% ============================================================================

\section{Heading Sizes: the \texttt{headings} Option}
\label{sec:headings}

Standard \LaTeX{} heading sizes can feel oversized, especially when combined
with large base fonts. KOMA-Script's \texttt{headings} option scales all
section-level headings uniformly.

\subsection{Available Values}

\begin{description}
    \item[\texttt{headings=big}]
        Large headings, similar to the standard \LaTeX{} defaults.
    \item[\texttt{headings=normal}]
        Moderate headings; a good balance between prominence and text density.
    \item[\texttt{headings=small}]
        Smaller headings that do not dominate the page. (\emph{Used in this
        document.})
\end{description}

With \texttt{headings=small}, the section headings above (``KOMA-Script:
\texttt{scrartcl} in Practice'', ``Why KOMA-Script?'', etc.) are noticeably
less imposing than the \texttt{big} setting would produce.

\subsection{Heading Fonts}

KOMA-Script also separates \emph{heading font} from \emph{body font}.
You can change all headings to sans-serif in the preamble:

\begin{verbatim}
\setkomafont{sectioning}{\normalfont\sffamily\bfseries}
\end{verbatim}

Or change only specific levels:

\begin{verbatim}
\setkomafont{section}{\normalfont\Large\bfseries}
\setkomafont{subsection}{\normalfont\large\itshape}
\end{verbatim}

% ============================================================================
% SECTION 4: THE DIV OPTION AND TYPEAREA
% ============================================================================

\section{Typearea: the \texttt{DIV} Option}
\label{sec:div}

The most conceptually distinctive KOMA-Script feature is its typearea
algorithm, which derives margins from the paper dimensions rather than
requiring the user to specify them manually.

\subsection{What DIV Does}

The \texttt{DIV} value $n$ divides the page height and width each into $n$
equal strips. The text block height occupies $n-2$ strips; the text block
width occupies $n-2$ strips. Margins are distributed 1 strip on the binding
side and 2 on the outer side (classical golden-ratio rule).

\begin{center}
\begin{tabular}{rll}
\toprule
\textbf{DIV} & \textbf{Text area (A4)} & \textbf{Typical use} \\
\midrule
 9  & 149\,mm \texttimes{} 213\,mm & Comfortable reading; wide margins \\
10  & 161\,mm \texttimes{} 228\,mm & Standard academic documents \\
11  & 170\,mm \texttimes{} 238\,mm & Slightly tighter \\
12  & 177\,mm \texttimes{} 247\,mm & Used in this document \\
14  & 188\,mm \texttimes{} 260\,mm & Dense; minimal margins \\
\bottomrule
\end{tabular}
\end{center}

\subsection{Specifying DIV}

In the \verb|\documentclass| options:

\begin{verbatim}
\documentclass[DIV=12]{scrartcl}
\end{verbatim}

Or dynamically after setting the font:

\begin{verbatim}
\usepackage{lmodern}
\KOMAoption{DIV}{calc}   % recalculate after font change
\end{verbatim}

Setting \texttt{DIV=calc} instructs KOMA-Script to choose an optimal value
automatically based on the current font and paper size.

% ============================================================================
% SECTION 5: KOMA vs. STANDARD CLASS COMPARISON
% ============================================================================

\section{Side-by-Side Comparison}
\label{sec:comparison}

\begin{center}
\begin{tabular}{p{4.5cm}p{4.5cm}p{4.5cm}}
\toprule
\textbf{Feature} & \textbf{Standard \texttt{article}} & \textbf{KOMA \texttt{scrartcl}} \\
\midrule
Margins & Manual via \texttt{geometry} & Auto via \texttt{DIV} (or \texttt{geometry}) \\
Paragraph spacing & First-line indent & Configurable (\texttt{parskip=}) \\
Heading sizes & Fixed large sizes & Configurable (\texttt{headings=}) \\
Subject/Subtitle & Not built in & \verb|\subject{}|, \verb|\subtitle{}| \\
Heading fonts & Hard to change & \verb|\setkomafont{}| \\
Caption formatting & \texttt{caption} package & Built in \\
KOMA-Options interface & No & Yes (\verb|\KOMAoption{}|) \\
\bottomrule
\end{tabular}
\end{center}

% ============================================================================
% SECTION 6: QUICK REFERENCE
% ============================================================================

\section{Quick Reference}
\label{sec:reference}

\subsection{Class Options Summary}

\begin{verbatim}
\documentclass[
    12pt,             % base font: 10pt, 11pt, 12pt
    a4paper,          % paper: a4paper, letterpaper, a5paper ...
    parskip=half,     % paragraph spacing: false, half, full
    headings=small,   % heading size: big, normal, small
    DIV=12,           % typearea divisor: integer or calc
    twoside,          % twoside / oneside
    open=right        % new chapters: right, left, any
]{scrartcl}
\end{verbatim}

\subsection{Useful KOMA-Script Commands}

\begin{description}
    \item[\texttt{\textbackslash KOMAoption\{option\}\{value\}}]
        Change a KOMA option inside the document body.
    \item[\texttt{\textbackslash setkomafont\{element\}\{font decl\}}]
        Set the font for a named KOMA element.
    \item[\texttt{\textbackslash addtokomafont\{element\}\{font decl\}}]
        Append to an existing KOMA font setting.
    \item[\texttt{\textbackslash subject\{...\}}]
        Short descriptive line above the title.
    \item[\texttt{\textbackslash subtitle\{...\}}]
        Subtitle below the main title.
    \item[\texttt{\textbackslash publishers\{...\}}]
        Publisher information on the title page.
\end{description}

\end{document}

% ============================================================================
% COMPILATION NOTES
% ============================================================================
% pdflatex article_vs_koma.tex
%
% If you use \addbibresource + \printbibliography, also run:
%   biber article_vs_koma
% then pdflatex twice more.
% ============================================================================
